\documentclass{article}
\usepackage[T1]{fontenc}
\usepackage{lmodern}
\usepackage{graphicx}
\usepackage{amsmath,amssymb}
\usepackage[a4paper,margin=2cm]{geometry}
\usepackage[absolute,overlay]{textpos}
\setlength{\TPHorizModule}{1mm}
\setlength{\TPVertModule}{1mm}
\usepackage[indonesian]{babel}
\usepackage{hyperref}
\setlength{\emergencystretch}{2em}
% unit: Halaman 42

\begin{document}

% === Halaman 42 (python/native) ===

Contoh: Caesar cipher   (ingat kembali materi di dalam kuliah Matdis ☺ )

• M = \{ huruf-huruf alfabet \}

• K = \{ k | k  adalah bilangan bulat dan 0 ≤ k ≤ 25 \}

• E = \{ Ek | k  K dan untuk semua plainteks m,




Ek(m) = (m + k) mod 26 \}

• D = \{ Dk | k  K dan untuk semua cipherteks c,




Dk(c) = (c – k) mod 26 \}

• C = M

42

\end{document}
